\chapter{Variational Calculus and Classical Mechanics}

\section{From functions to functionals}

A function takes a number and returns a number.  A functional takes a function and returns a number.  The action of a particle moving along \(q(t)\) is
\[
  S[q]=\int_{t_i}^{t_f}L(q,\dot q,t)\dd t.
\]
The physical path is the one for which the action is stationary under small changes of the path with fixed endpoints.

Let
\[
  q(t)\to q(t)+\epsilon\eta(t),
  \qquad
  \eta(t_i)=\eta(t_f)=0.
\]
The first-order change is
\[
  \delta S
  =
  \left.\frac{\dd}{\dd\epsilon}S[q+\epsilon\eta]\right|_{\epsilon=0}.
\]
Using the chain rule,
\[
  \delta S
  =
  \int_{t_i}^{t_f}
  \left(
  \frac{\p L}{\p q}\eta
  +
  \frac{\p L}{\p\dot q}\dot\eta
  \right)\dd t.
\]
Integrate the second term by parts:
\[
  \int_{t_i}^{t_f}\frac{\p L}{\p\dot q}\dot\eta\,\dd t
  =
  \left[\frac{\p L}{\p\dot q}\eta\right]_{t_i}^{t_f}
  -
  \int_{t_i}^{t_f}
  \frac{\dd}{\dd t}\left(\frac{\p L}{\p\dot q}\right)\eta\,\dd t.
\]
The boundary term vanishes.  Therefore
\[
  \delta S
  =
  \int_{t_i}^{t_f}
  \left[
  \frac{\p L}{\p q}
  -
  \frac{\dd}{\dd t}\left(\frac{\p L}{\p\dot q}\right)
  \right]\eta(t)\dd t.
\]
For this to vanish for every \(\eta(t)\), the bracket must vanish:
\[
  \frac{\dd}{\dd t}\left(\frac{\p L}{\p\dot q}\right)-\frac{\p L}{\p q}=0.
\]
This is the Euler--Lagrange equation.

\begin{example}[Newton's law]
Take
\[
  L=\frac12m\dot q^2-V(q).
\]
Then
\[
  \frac{\p L}{\p\dot q}=m\dot q,\qquad
  \frac{\p L}{\p q}=-\frac{\dd V}{\dd q}.
\]
The Euler--Lagrange equation gives
\[
  m\ddot q=-\frac{\dd V}{\dd q},
\]
which is Newton's second law with force \(F=-V'\).
\end{example}

\section{Many coordinates}

For coordinates \(q_1,\ldots,q_n\),
\[
  S[q]=\int L(q_1,\ldots,q_n,\dot q_1,\ldots,\dot q_n,t)\dd t.
\]
Varying one coordinate at a time gives
\[
  \frac{\dd}{\dd t}\left(\frac{\p L}{\p\dot q_i}\right)
  -
  \frac{\p L}{\p q_i}=0,
  \qquad i=1,\ldots,n.
\]
The canonical momentum is
\[
  p_i=\frac{\p L}{\p\dot q_i}.
\]
When the velocities can be solved in terms of \(q\) and \(p\), define the Hamiltonian
\[
  H(q,p,t)=\sum_i p_i\dot q_i-L.
\]
Differentiating,
\[
  \dd H
  =
  \sum_i \dot q_i\,\dd p_i
  +\sum_i p_i\,\dd\dot q_i
  -\sum_i\frac{\p L}{\p q_i}\dd q_i
  -\sum_i\frac{\p L}{\p\dot q_i}\dd\dot q_i
  -\frac{\p L}{\p t}\dd t.
\]
The \(p_i\,\dd\dot q_i\) terms cancel because \(p_i=\p L/\p\dot q_i\), and the equations of motion give \(\p L/\p q_i=\dot p_i\).  Hence
\[
  \dd H=\sum_i\dot q_i\,\dd p_i-\sum_i\dot p_i\,\dd q_i-\frac{\p L}{\p t}\dd t.
\]
Therefore
\[
  \dot q_i=\frac{\p H}{\p p_i},
  \qquad
  \dot p_i=-\frac{\p H}{\p q_i}.
\]

\section{The harmonic oscillator}

For
\[
  L=\frac12m\dot q^2-\frac12m\omega^2q^2,
\]
the equation of motion is
\[
  \ddot q+\omega^2q=0.
\]
The general solution is
\[
  q(t)=A\cos\omega t+B\sin\omega t.
\]
The momentum is \(p=m\dot q\), and the Hamiltonian is
\[
  H=\frac{p^2}{2m}+\frac12m\omega^2q^2.
\]
The oscillator is the seed of free quantum field theory.  A free field is an infinite collection of oscillators, one for each wave vector.

\section{Noether's theorem for mechanics}

Suppose the action is unchanged under a continuous transformation
\[
  q_i(t)\to q_i(t)+\epsilon\Delta q_i(t).
\]
If the Lagrangian changes only by a total derivative,
\[
  \delta L=\epsilon\frac{\dd F}{\dd t},
\]
then there is a conserved quantity.  Directly,
\[
  \delta L=
  \sum_i\frac{\p L}{\p q_i}\Delta q_i
  +
  \sum_i\frac{\p L}{\p\dot q_i}\frac{\dd}{\dd t}\Delta q_i.
\]
Use the equations of motion \(\p L/\p q_i=\dot p_i\):
\[
  \delta L=
  \sum_i\dot p_i\Delta q_i+\sum_i p_i\frac{\dd}{\dd t}\Delta q_i
  =
  \frac{\dd}{\dd t}\left(\sum_i p_i\Delta q_i\right).
\]
Comparing with \(\delta L=\dd F/\dd t\),
\[
  \frac{\dd}{\dd t}
  \left(
  \sum_i p_i\Delta q_i-F
  \right)=0.
\]
The conserved charge is
\[
  Q=\sum_i p_i\Delta q_i-F.
\]

\begin{example}[Time translation and energy]
If \(L\) has no explicit time dependence, a shift \(t\to t+\epsilon\) is a symmetry.  The associated conserved quantity is the Hamiltonian \(H\).  This can also be seen directly:
\[
  \frac{\dd H}{\dd t}
  =
  \frac{\p H}{\p q}\dot q+\frac{\p H}{\p p}\dot p+\frac{\p H}{\p t}
  =
  -\dot p\,\dot q+\dot q\,\dot p+\frac{\p H}{\p t}.
\]
If \(\p H/\p t=0\), then \(\dd H/\dd t=0\).
\end{example}

\section{Fields as infinitely many coordinates}

Discretize space into lattice points \(\bm x_n\).  A scalar field \(\phi(t,\bm x)\) becomes many coordinates
\[
  q_n(t)=\phi(t,\bm x_n).
\]
The continuum limit replaces sums by integrals:
\[
  \sum_n \Delta V\,(\cdots)\to \int\dd^3x\,(\cdots).
\]
A field action has the form
\[
  S[\phi]=\int\dd t\int\dd^3x\,\Lag(\phi,\p_t\phi,\nabla\phi).
\]
The field Euler--Lagrange equation is derived in the next chapter by exactly the same integration-by-parts idea.

\section*{Exercises}
\addcontentsline{toc}{section}{Exercises}

\begin{exercise}
Derive the equation of motion for
\[
  L=\frac12m\dot q^2-\frac{\lambda}{4}q^4.
\]
\begin{answer}
\(m\ddot q+\lambda q^3=0\).
\end{answer}
\end{exercise}

\begin{exercise}
For the free particle \(L=\frac12m\dot q^2\), use spatial translation \(q\to q+\epsilon\) to find the conserved Noether charge.
\begin{answer}
\(\Delta q=1\), \(F=0\), so \(Q=p=m\dot q\).
\end{answer}
\end{exercise}

